\documentclass[10pt,a4paper]{amsart}
\usepackage[margin=19mm]{geometry}
\usepackage{amsmath,amssymb,mathtools}
\usepackage[scr=rsfs,cal=euler]{mathalfa}
\usepackage{xfrac}
\usepackage[unicode,hidelinks,bookmarks=false]{hyperref}
\newcommand{\ie}{i.e.\ }
\newcommand{\cf}{cf.\ }
\newcommand{\resp}{respectively\ }
\newcommand{\Subsection}{\subsection}
\providecommand{\nobreakdash}{\nobreak\hbox{-}}
\setlength{\emergencystretch}{2em}
\multlinegap0pt
\usepackage{amssymb}
\usepackage{xcolor}
\usepackage[shortlabels]{enumitem}
\usepackage{tikz}
\usepackage{bm}
\usepackage{csquotes}
\usepackage{float}
\newtheorem{lemma}{Lemma}
\def\acts{\curvearrowright}
\DeclareMathOperator{\ppww}{pw}
\newcommand{\pw}{\tau_{\ppww}}
\newcommand{\pwom}{\tau_{\ppww}^{\spOm}}
\newcommand{\spOm}{\bm{\Omega}}
\title{All mixed identities are singular in groups with no algebraicity}
\author{Paolo Marimon, Michael Pinsker}
\date{Source: arXiv:2606.24741v2 (CC BY 4.0)}
\begin{document}
\maketitle

\noindent\emph{Source and licence.} All mixed identities are singular in groups with no algebraicity. Version arXiv:2606.24741v2 (CC BY 4.0), distributed by arXiv under the Creative Commons Attribution 4.0 International licence, http://creativecommons.org/licenses/by/4.0/, as recorded in the arXiv licence metadata for this exact version and shown on the abstract page https://arxiv.org/abs/2606.24741v2. Full licence notice: \texttt{licenses/arxiv-2606.24741-CC-BY-4.0.txt}.\par\smallskip

\noindent\textit{Editorial scope.} Counted unit: the article's lemma technical (source0.tex lines 1699--1710). The article's complete original proof of it is delivered with it (source0.tex lines 1711--1734, one \texttt{proof} environment of 24 lines). No mathematics is written, repaired or re-derived: the statement and the proof are the article's own lines, in the article's own notation, and no retained line is substituted or re-flowed.  No further source-local context is needed: every cross-reference of every retained line resolves inside the two delivered spans.  Nothing was omitted from any retained span, so the package carries no omission record.  No retained line contains a citation, so the wrapper carries no bibliography and no bibliography entry.  No retained line contains a figure environment, an image-inclusion command or drawing code, so no image is delivered and the package declares no asset dependency.  Editorial formatting only, in addition: an AMS \texttt{amsart} preamble that restates the article's own declarations and macros used by the retained lines, namely the environments \texttt{lemma} on separate counters, together with the standard packages those lines need.  The retained environments print in this document as Lemma 1 in order of appearance, whereas the article prints the same items with the numbering of its own full document.  The complete native source of the article as distributed in the arXiv e-print for this exact version is bundled unchanged as \texttt{sources/203546/source0.tex}.
\begin{lemma}\label{lemma:slightly technical} 
Let $G\acts\Omega$ be any group action. 
%{Let $\spOm$ be a %Hausdorff
%topological space and let $G\acts\Omega$ be a group action by homeomorphisms of $\spOm$.} 
Suppose that for any regular word ${w}\in G* F_r$ and for any finite $B\subseteq \Omega$, there are $p_1, \dots, p_r\in G_B$ such that ${w}(p_1, \dots, p_r)\neq 1$. Then, for any regular word $w\in G* F_r$, the set 
\begin{equation}\label{eq:thesolutionset}
   \{ (g_1, \dots, g_r)\in G^r\ \vert\ w(g_1, \dots, g_r){\neq} 1\}\; 
\end{equation}
{is dense in $G^r$ with respect to $\pw$.}
%the topology of pointwise convergence $\pwom$}  {induced by the action.}
% is nowhere dense in $G^r$ with respect to the topology of pointwise convergence. 
\end{lemma}

\begin{proof} {Let $\mathcal{U}$ be a non-empty basic open subset of $G^r$. We may write}
\[
  \mathcal{U}
  = \prod_{i=1}^r G_{(\overline{b}_i, \overline{b}'_i)},\]
{where $\overline{b}_i$ and $\overline{b}_i'$ are tuples from $\Omega$.} {To prove density,} it suffices to find $(g_1,\ldots,g_r)\in \mathcal U$ with $w(g_1,\ldots,g_r)\neq 1$. Choose
$\mathbf f=(f_1,\ldots,f_r)\in\mathcal U$, and let $B\subseteq \Omega$ be the
finite set consisting of all entries of the tuples
$\bar b_1,\ldots,\bar b_r$. Let $w_{\mathbf f}$ be the word obtained by {replacing} each occurrence of $x_i$ by $f_ix_i$, for all {$i\in\{1, \dots, r\}$.} Clearly, $w_{\mathbf f}$ is regular, and for every
$\mathbf p=(p_1,\ldots,p_r)\in G^r$, 
\begin{equation}\label{eq:substitution-evaluation}
  w_{\mathbf f}(p_1,\ldots,p_r)
  =w(f_1p_1,\ldots,f_rp_r).
\end{equation}
By assumption, there are
$p_1,\ldots,p_r\in G_B$ such that $w_{\mathbf f}(p_1,\ldots,p_r)\neq 1$.
For  $1\leq i\leq r$, set $g_i:=f_i p_i$. Since every $p_i$ fixes $B$ pointwise, $g_i$ agrees with $f_i$ on
$\bar b_i$, and so $(g_1,\ldots,g_r)\in\mathcal U$. By \eqref{eq:substitution-evaluation}, $w(g_1, \dots, g_r)\neq1$.
%Choose $a_0\in X$ such that $w(g_1, \dots, g_r)(a_0)\neq a_0$ and let $\overline{a}:=(a_0, \dots, a_{n+1})$ be the  the sequence for $w$ evaluated at $(g_1, \dots, g_r)$. Consider the following basic open neighbourhood of $\mathbf g$:
%  \[\mathcal W
%  :=\mathcal U\cap\prod_{i=1}^r G_{(\overline{a}, g_i(\overline{a}))}\;.\]
%By construction for any  $(h_1,\ldots,h_r)\in\mathcal W$, $w(h_1, \dots, h_r)(a_0)\neq a_0$. In particular, $\mathcal{W}\cap\mathfrak{S}=\emptyset$. Thus every non-empty basic
%open set contains a non-empty open subset disjoint from $\mathfrak{S}$. Equivalently,
%$\mathfrak{S}$ is nowhere dense.
\end{proof}

\end{document}
